\chapter{Background}

The novel coronavirus (COVID-19) pandemic has emerged as one of the most significant global health crises of the 21st century, affecting nations across the world irrespective of economic status or healthcare capacity. Originating in Wuhan, China, in late 2019, the virus has demonstrated the profound interconnectedness of modern societies, exposing vulnerabilities in public health systems, economies, and the fabric of global cooperation \parencite{Verma2021AReturns, Shrestha2020TheGlobalization, LegeseFeyisa2020TheReview, MaitalEllaBarzani2020EResearch}.

\section{The Unprecedented Impact of COVID-19}

COVID-19 has underscored the fragility of the global economy, causing widespread economic disruptions, from the decimation of industries such as travel and tourism to unprecedented levels of unemployment and poverty \parencite{Verma2021AReturns, LegeseFeyisa2020TheReview}. The pandemic has catalyzed a deep recession, with the World Bank and the International Monetary Fund predicting a contraction in global GDP unparalleled in recent history \parencite{Verma2021AReturns, Shrestha2020TheGlobalization, LegeseFeyisa2020TheReview, MaitalEllaBarzani2020EResearch}. The abrupt halt in economic activity has led to a significant decline in stock markets worldwide, reflecting investor uncertainty and fear \parencite{Verma2021AReturns}.

The response to the pandemic has varied globally, with nations implementing lock-downs, social distancing measures, and mass testing in attempts to curb the spread of the virus. These measures, while necessary for public health, have further strained economies, leading to a comprehensive debate on balancing health outcomes with economic sustainability \parencite{Shrestha2020TheGlobalization, LegeseFeyisa2020TheReview}.

\section{Learning from History for Future Preparedness}

Historically, pandemics have been recurrent challenges for humanity, with the 1918 Spanish Flu being a notable reference point for COVID-19 due to its widespread impact and high mortality rate. Similar to past pandemics, COVID-19 has highlighted the critical need for robust healthcare systems, rapid response mechanisms, and global cooperation in surveillance and vaccine development \parencite{Shrestha2020TheGlobalization, MaitalEllaBarzani2020EResearch}. The pandemic has also emphasized the importance of non-pharmaceutical interventions, such as social distancing and hygiene practices, which have become the first line of defense in the absence of a vaccine \parencite{LegeseFeyisa2020TheReview, MaitalEllaBarzani2020EResearch}.

\section{The Role of Technology and Innovation}

The COVID-19 crisis has accelerated the adoption of technology and innovation across various sectors. Telemedicine, remote work, and e-learning have become the new norm, suggesting a long-term transformation in how societies operate. This shift presents an opportunity to re-imagine future responses to pandemics, with a focus on digital infrastructure, data analytics for disease surveillance, and a more agile healthcare delivery model \parencite{Shrestha2020TheGlobalization, MaitalEllaBarzani2020EResearch}.

\section{A Call for Global Solidarity and Preparedness}

The COVID-19 pandemic has served as a stark reminder of the global community's interconnectedness and the collective vulnerability to emerging pathogens. It underscores the necessity for a unified global response strategy, investments in healthcare infrastructure, and the importance of preparedness plans that are adaptable and scalable in the face of future pandemics. The lessons learned from COVID-19 should guide international policy, research, and cooperation, ensuring that the world is better equipped to face similar challenges in the future \parencite{Verma2021AReturns, Shrestha2020TheGlobalization, LegeseFeyisa2020TheReview, MaitalEllaBarzani2020EResearch}.

\section{Supply Chain Management and the Needs for Constant Improvement}

Supply chain management is a critical aspect of any industry, with the food industry being particularly sensitive due to the nature of its products and the importance of maintaining quality and safety. In today’s rapidly evolving global market, various vulnerabilities within the supply chain can have significant impacts on both businesses and consumers. One of the pressing concerns is identifying specific vulnerabilities within supply chains. For instance, what happens if we encountered health-related problems in the food industry due to supply chain issues? Understanding these vulnerabilities is crucial for developing robust strategies to mitigate risks. 

If and when we encounter a real-world example of a health-related problem in the food industry. How was it caused, and what were the outcomes? These situations highlight the importance of having complete end-to-end visibility in the supply chain. However, achieving such transparency comes with its own set of challenges \parencite{Wagner2006AnVulnerability, Elleuch2016ResilienceReview}. To address these vulnerabilities, avoiding over-reliance on single sources is a key strategy in mitigating risk. What are the strategies we can use to diversify our supply sources effectively? Furthermore, how adaptable and flexible the current supply chain is in responding to fluctuations in demand or supply \parencite{Yin2022SupplyFsQCA, Wang2024TheCOVID-19, Lin2021TheManufacturers}. Balancing cost efficiency with other factors like safety and quality is another challenge, especially in the context of global operations such as temperature-controlled transport. What is the optimal inventory strategy to achieve this balance? How does this strategy apply on a global scale?

Collaboration with suppliers and partners is essential for improving transparency, trust, and joint risk management. During unprecedented events like pandemics, what kinds of collaboration have proven effective? Similarly, how can technologies like AI and IoT help solve tracking and traceability challenges, thereby enhancing the overall supply chain \parencite{Wang2024TheCOVID-19, Simatupang2002TheChain, Singh2023Post-COVIDVaccination, Ramanathan2014SupplyPartnerships, Chen2017SupplyAgenda}? With increasing scrutiny of regulating and governing bodies such as EU-Parliament on environmental impacts and green supply chain \parencite{Amann2014DrivingUnion, Moazzem2022EnvironmentalProducts}, ensuring compliance with global regulations while maintaining a commitment to sustainability is another critical consideration. 

What strategies can be employed to reduce waste and ensure the proper use of resources? Moreover, do you have comprehensive crisis management plans that address supply chain disruptions, and how were these plans tested during the COVID-19 pandemic? Rapid recovery from supply chain disruptions requires well-thought-out strategies. Can you provide examples of successful recovery strategies? During the COVID-19 pandemic, what were the most significant challenges faced, especially concerning perishable goods? Looking ahead, what are the main concerns regarding the potential impact of future pandemics on perishable goods, and what measures can the industry take to improve resilience \parencite{Chowdhury2021COVID-19Review, Paul2021SupplyPandemic, Ivanov2017LiteratureChain, Chen2019BuildingIndustry}? 

In conclusion, the resilience and adaptability of the supply chain are of paramount importance, especially in the face of unforeseen challenges. By exploring these questions, we can better understand the intricacies of supply chain management and develop strategies to ensure its robustness in the future.

\subsection{Identifying Research Gaps}

The studies reviewed thus far provide a foundational understanding of how supply chains have adapted during the COVID-19 pandemic, particularly in handling perishable goods. However, these insights reveal several areas where further investigation is crucial. There remains a significant gap in understanding the full impact of technological advancements on supply chain resilience. Emerging technologies such as artificial intelligence (AI) and blockchain have shown potential to revolutionize supply chain management by enhancing efficiency, transparency, and responsiveness. Detailed research is needed to evaluate how these technologies can be specifically leveraged to bolster the resilience of supply chains managing perishable goods. This is critical for ensuring that food supplies remain safe and efficient, minimizing waste and improving distribution even under disruptive conditions. 

Additionally, the adaptations and strategies developed for specific sectors during the pandemic, such as the implementation of advanced digital solutions in regulated environments like healthcare, suggest that similar strategies could be applied to the perishable goods sector. More focused studies are required to explore how these strategies can be tailored to meet the unique demands and challenges of different industries, ensuring that supply chains can withstand not only current disruptions but also future crises.

Furthermore, while the integration of online and offline channels in retail has been examined, primarily through models like 'buy online and pick up in store' (BOPS), there is a need to delve deeper into the scalability and practical implementation of such models across various supply chain networks. The mathematical models provided by other studies have started to address these issues, but more comprehensive research is necessary to establish best practices and identify the optimal conditions under which these integrated strategies enhance supply chain resilience. 

This includes investigating the sector-specific challenges and opportunities that arise when adapting supply chain designs to efficiently manage perishable goods. Each sector presents unique conditions and regulatory frameworks that can affect the implementation of resilient supply chain strategies. Therefore, research aimed at uncovering these nuances and developing sector-specific adaptations could provide invaluable guidance for businesses seeking to enhance their operational resilience in the face of ongoing and future disruptions.

% The COVID-19 pandemic has brought about unprecedented changes in consumer behavior and retail dynamics, presenting a compelling need for in-depth analysis. This research is motivated by several key factors:

% \begin{itemize}
%     \item \textbf{Transformation in Consumer Behavior}: The pandemic has radically altered how consumers shop, work, and live. New personal circumstances, such as changes in income and leisure time, have significantly influenced consumer attitudes and behaviors. Consumers have become more conscious of environmental, health, and cost factors, leading to a preference for locally-sourced products and neighborhood stores. Notably, there has been a significant rise in digital commerce, especially among new or low-frequency users, which is expected to continue post-pandemic \parencite{standish_2020_covid19}.

%     \item \textbf{Lasting Impact on Retail and E-Commerce}: The shifts in consumer behavior are not just transient reactions to the pandemic but are likely to have long-term implications for the retail industry. Companies have an opportunity to help shape the 'next normal' as many of the longer-term changes in consumer behavior are still forming. This presents a unique challenge and opportunity for retailers and consumer packaged goods companies to adapt and evolve in response to these changes \parencite{fabius_2020_how}.

%     \item \textbf{Need for Strategic Adaptation}: Understanding these behavioral shifts is crucial for businesses to develop effective strategies for the post-pandemic market. This includes rethinking supply chain management, marketing strategies, and logistics operations to align with the evolving consumer preferences and shopping habits.

%     \item \textbf{Resilience and Response to Future Uncertainties}: Analyzing how the pandemic has transformed consumer behavior and retail dynamics can provide insights into how businesses and consumers respond to uncertainties. This knowledge can be instrumental in preparing for future crises, ensuring more resilient and adaptable business models.
    
% \end{itemize}

% this research is vital for comprehending the extensive and potentially lasting changes brought about by the COVID-19 pandemic in the realm of consumer behavior and retail. It offers valuable insights for businesses to adapt to the evolving market and consumer needs in the post-pandemic era, and to prepare for future challenges. This is a citation \parencite[p.~12]{guo2022has}. \textcite[p.~12]{guo2022has} has proposed something. In Figure \ref{fig:covid_art2} there is covid art.

% \begin{figure}[htbp] % Positioning option: h=here, t=top, b=bottom, p=page of floats
%   \centering % Centers the image
%   \fbox{\includegraphics[width=0.8\textwidth]{figure/covid_art.jpg}}
%   \caption{Sample COVID art.}
%   \label{fig:covid_art} 
% \end{figure}

% \begin{figure}[htbp] % Positioning option: h=here, t=top, b=bottom, p=page of floats
%   \centering % Centers the image
%   \fbox{\includegraphics[width=0.8\textwidth]{figure/art.jpg}}
%   \caption{Some caption.}
%   \label{fig:covid_art2} 
% \end{figure}

% \parencite{guo2022has}
% \textcite{guo2022has}
% \textcite[p.~12]{guo2022has}

% \begin{itemize}
%     \item Transformation in Consumer Behavior: The pandemic has radically altered how consumers shop, work, and live. New personal circumstances, such as changes in income and leisure time, have significantly influenced consumer attitudes and behaviors. Consumers have become more conscious of environmental, health, and cost factors, leading to a preference for locally-sourced products and neighborhood stores. Notably, there has been a significant rise in digital commerce, especially among new or low-frequency users, which is expected to continue post-pandemic (Standish, 2020).
    
%     \item Lasting Impact on Retail and E-Commerce: The shifts in consumer behavior are not just transient reactions to the pandemic but are likely to have long-term implications for the retail industry. Companies have an opportunity to help shape the 'next normal' as many of the longer-term changes in consumer behavior are still forming. This presents a unique challenge and opportunity for retailers and consumer packaged goods companies to adapt and evolve in response to these changes (Fabius, 2020).
    
%     \item Need for Strategic Adaptation: Understanding these behavioral shifts is crucial for businesses to develop effective strategies for the post-pandemic market. This includes rethinking supply chain management, marketing strategies, and logistics operations to align with the evolving consumer preferences and shopping habits.
% \end{itemize}
